\section{Literature Review and Related Work}

\subsection{Statistical, Machine-Learning, and Deep Forecasters}
Forecasting methods differ in how they trade inductive bias, data demand, and
computational cost. Autoregressive and seasonal models remain informative
baselines because their failure modes are easy to inspect. Tree ensembles such
as XGBoost and LightGBM extend this setting to nonlinear lag, rolling, weather,
and calendar interactions without requiring a sequence encoder
\cite{chen2016xgboost,ke2017lightgbm}. Recurrent LSTM and GRU models instead
learn stateful temporal representations \cite{hochreiter1997lstm,cho2014gru}.
Large benchmark archives, including the Monash Forecasting Repository, further
support reproducible comparison across heterogeneous time-series domains
\cite{godahewa2021monash}.
Transformer families emphasize longer dependencies: Informer reduces attention
cost, Autoformer and FEDformer use decomposition, Crossformer and iTransformer
model cross-variable structure, and TimesNet reorganizes temporal variation
into two-dimensional patterns
\cite{zhou2021informer,wu2021autoformer,zhou2022fedformer,zhang2023crossformer,liu2024itransformer,wu2023timesnet}.
These gains do not imply that greater architectural complexity is always
necessary. DLinear and PatchTST show that simple decomposition or patch-based
inductive biases can remain highly competitive \cite{zeng2023dlinear,nie2023patchtst},
whereas TimeMixer uses explicit multiscale decomposition and mixing
\cite{wang2024timemixer}. Event-TimeRAF therefore retains seasonal and boosted
tree baselines instead of treating a deep backbone as the default winner.

\subsection{Time-Series Foundation Models}
Time-series foundation models (TSFMs) seek transfer across datasets rather than
training a separate representation from scratch for each series. TimesFM uses a
decoder-only forecasting design; Moirai trains a universal forecasting
Transformer; MOMENT supports several time-series tasks; Chronos represents
values as tokens; and Tiny Time Mixers emphasizes compact zero- and few-shot
transfer
\cite{das2024timesfm,woo2024moirai,goswami2024moment,ansari2024chronos,ekambaram2024ttm}.
Timer, Timer-XL, UniTime, and Time-LLM explore generative pretraining,
long-context scaling, cross-domain unification, or language-model reprogramming
\cite{liu2024timer,liu2025timerxl,liu2024unitime,jin2024timellm}. Frozen
pretrained models have also been evaluated as general computation engines for
time-series analysis \cite{zhou2023onefitsall}. Collectively, these studies
motivate zero-shot evaluation, but they also show that transfer performance
depends on input representation and domain shift. The present study therefore
compares frozen Chronos-Bolt against locally supervised models on identical
forecast origins rather than assuming foundation-model superiority.

\subsection{Retrieval-Augmented Forecasting}
Retrieval augmentation separates parametric knowledge from an external memory.
In natural-language processing, RAG and Dense Passage Retrieval established the
retrieve-then-condition pattern \cite{lewis2020rag,karpukhin2020dpr}. In time
series, retrieval has been studied through diffusion forecasting and assisted
generation workflows \cite{liu2024ratd,ang2024tsgassist}. TimeRAF provides the
closest methodological reference: it constructs a time-series knowledge base,
trains an MLP dual-encoder retriever, selects top-$k$ candidates, and integrates
candidate embeddings into a frozen TSFM through Channel Prompting
\cite{zhang2025timeraf}. Its ablations separate retrieval quality, integration,
candidate count, knowledge-base composition, and backbone choice.

Event-TimeRAF retains the external-memory principle but addresses a different
experimental question. It uses transparent random, cosine, and event-context
similarity instead of a learned dual encoder, and it uses feature- or
forecast-level fusion instead of internal Channel Prompting. This distinction
prevents the current implementation from being described as a reproduction of
TimeRAF; it is an audited environmental adaptation designed to test whether
contextual historical analogues help a local PM2.5 task.

\subsection{Air-Quality Forecasting and Official Context}
PM2.5 forecasting commonly combines pollutant history with temperature,
humidity, wind, precipitation, and pressure. Recent reviews describe a shift
from single-station recurrent models toward attention-based, hybrid, and
spatiotemporal systems \cite{zhou2024pm25review,zhang2023sparsepm25}. AirFormer,
for example, separates deterministic spatiotemporal modeling from stochastic
uncertainty estimation at nationwide scale \cite{liang2023airformer}. Such
multi-station designs address spatial dependence, whereas the present study
examines a county-level aggregate and the provenance of contextual evidence.
EPA AirData supplies regulated PM2.5 observations, NOAA Integrated Surface
Database supplies hourly meteorology, and NOAA Storm Events supplies structured
hazard records \cite{epaairdata2026,noaaisd2026,noaastormevents2026}. NOAA
Hazard Mapping System products could add smoke evidence but were not included in
the reported run \cite{noaahms2026}.

\subsection{Distribution Shift and Forecast Explanations}
Temporal nonstationarity can alter means, variances, and relationships between
inputs and future values. Adaptive normalization methods address this problem by
estimating statistics at local temporal scales \cite{liu2023san}; Event-TimeRAF
instead uses interpretable shift indicators to audit regimes without claiming
that shift has been eliminated. Explanation methods also differ in scope. SHAP
attributes predictions to model inputs \cite{lundberg2017shap}, while
information-theoretic saliency supports counterfactual inspection of
probabilistic forecasts \cite{raman2023forecast}. The implemented explanation
layer is narrower: it records signed XGBoost effects, retrieved cases, event
identifiers, drift components, and an uncertainty proxy. These records support
traceability, but they do not establish causal effects of weather or events.
